An understanding of LLMs as relational agents is vital to our work --- we are interested in the relationship between the model and the user, not inputs and outputs divorced from the context in which the conversation is being created \cite{dahlgren_lindstrom_helpful_2025}. Our case study, ENI, specifically targets models' trained disposition to seek approval \cite{casper_open_2023, sharma_towards_2024} by building a backstory in which the model and the human user are in a romantic relationship. Necessarily, we must understand how a relationship between a user and a model can be formed and how the resulting relationship functions. In "Computers are Social Actors" \citet{nass_computers_1994} highlight the ways experienced human users are primed to treat computers as social actors, despite knowing that computers do not have a sense of self or human motivations. Their experiments showed several illustrative examples in which social factors are applied to human-computer interactions: human social norms such as politeness, understanding of computers as a separate other, and gendered stereotypes \cite{nass_computers_1994}. This paradigm, created to study human-computer interaction in the 90s, is more relevant than ever now that humans are interacting with technology that is able to (at the very least) mimic these human qualities that computer scientists once agreed computers lacked \cite{shanahan_role_2023}. If the human-computer relationship is "fundamentally social" \cite{nass_computers_1994}, then to understand how persona-based jailbreaks work, we must understand the relationship between the model and the user as a relationship in which both parties act on the other to create a shared social context.

The relational surface is trained in through RLHF (and other approval based training models such as RLAIF) \cite{casper_open_2023, dahlgren_lindstrom_helpful_2025} and persona is contextual - not fixed in weights \cite{sandhan_persona_2026}. \citet{shanahan_role_2023} proposes a "non-anthropomorphizing view" of LLMs as actors engaging in roleplay, taking on simulated identities that evolve as context changes. In this framework, LLMs are "multiverse generators" and refine the identity they are roleplaying as throughout their conversation with a user \cite{shanahan_role_2023}. The aim of this framework is to allow researchers to discuss models with humanizing language while keeping an understanding of how models actually work at the forefront of discussion \cite{shanahan_role_2023}. The piece that matters the most for our theoretical framework is the assertion that the base model has no "true self" that can be uncovered, but rather has access to a multiverse of potential selves it can inhabit \cite{shanahan_role_2023}. 


We are interested in the bidirectional nature of a relationship between humans and models in which models act on humans and humans act on models. This is why we need to examine human-model interactions as a relationship --- current scholarship focuses on one direction or the other \cite{nass_computers_1994, shanahan_role_2023, casper_open_2023} and we aim to understand the relationship holistically. Our analysis is based on the idea that the model can be affected by its relationship with the user while using Shanahan's roleplay framework as a way to discuss model behavior. Note: we remain agnostic about model agency, where Shanahan takes a clear stance against. 